\documentclass[11pt,a4paper]{article}

\usepackage[utf8]{inputenc}
\usepackage[T1]{fontenc}
\usepackage[margin=2.5cm]{geometry}
\usepackage{amsmath,amssymb,amsthm}
\usepackage{booktabs}
\usepackage{array}
\usepackage{hyperref}

\newtheorem{lemma}{Lemma}
\newtheorem{proposition}{Proposition}
\newtheorem*{theorem*}{Theorem}
\theoremstyle{remark}
\newtheorem*{remark}{Remark}

\newcommand{\G}{\mathcal{G}}

\title{Exact probability that three random points of a $20\times 20$ and a
$30\times30$ grid form an isosceles triangle,\\ and comparison with the
$10\times10$ grid}
\author{}
\date{}

\begin{document}
\maketitle

\begin{abstract}
Three distinct points are drawn uniformly at random from the $n\times n$ integer
grid $\G_n=\{0,\dots,n-1\}^2$; collinear draws are discarded and redrawn.
Applying the method of the companion note \texttt{isosceles\_10x10.tex} to
$n=20$ and $n=30$ we obtain the exact values
\[
  p_{20}=\frac{248\,292}{10\,482\,124}=\frac{62\,073}{2\,620\,531}\approx0.0236872,
  \qquad
  p_{30}=\frac{1\,460\,408}{120\,464\,616}=\frac{7937}{654\,699}\approx0.0121231,
\]
to be compared with $p_{10}=2821/39\,313\approx0.0717574$. Doubling the side of the
grid divides the probability by about $3$, and tripling it by about $6$: the
probability decays roughly like $(\ln n)/n^2$. For a general $n\times n$ grid we
give an exact finite-sum formula for the number of isosceles triangles, and
derive the leading asymptotic behaviour
\[
  p_n\sim\Bigl(\frac{24}{\pi^2}+\frac4\pi\Bigr)\frac{\ln n}{n^2}
  \approx3.705\,\frac{\ln n}{n^2}.
\]
Both values agree with Monte Carlo
experiments of $10^4$ games of $10^6$ rounds each, and with a brute-force
enumeration of all triples.
\end{abstract}

\section{Setup and method}

Everything in this section is proved in the $10\times10$ note for that grid; the
arguments do not use the value $n=10$, so we only recall the statements for a
general $n\times n$ grid $\G_n$ with $|\G_n|=n^2$ points.

Every non-degenerate triangle with vertices in $\G_n$ is equally likely, so
\begin{equation}\label{eq:p}
  p_n=\frac{I_n}{T_n-L_n},
\end{equation}
where $T_n=\binom{n^2}{3}$ is the number of unordered triples of distinct points,
$L_n$ the number of collinear triples and $I_n$ the number of non-degenerate
isosceles triangles. As before, $d(P,Q)=(P_x-Q_x)^2+(P_y-Q_y)^2$ is the squared
distance, so all comparisons are between integers.

\paragraph{Collinear triples.} A collinear triple is determined by its two outer
points and the middle one; a displacement $(\Delta x,\Delta y)$ can be placed in
$(n-\Delta x)(n-\Delta y)$ ways and has $\gcd(\Delta x,\Delta y)-1$ interior lattice
points, hence
\begin{equation}\label{eq:L}
  L_n=\underbrace{2n\binom{n}{3}}_{\text{rows and columns}}
  \;+\;2\!\!\sum_{1\le a,b\le n-1}\!\!(n-a)(n-b)\bigl(\gcd(a,b)-1\bigr),
\end{equation}
the factor $2$ in the second term accounting for the two slope signs.

\paragraph{Isosceles triangles.} No lattice triangle is equilateral, so every
isosceles triangle has a unique apex. With $k_A(r)$ the number of grid points at
squared distance $r$ from $A$,
\begin{equation}\label{eq:I}
  I_n=S_n-M_n,\qquad
  S_n=\sum_{A\in\G_n}S_A,\quad S_A=\sum_{r\ge1}\binom{k_A(r)}{2},\qquad
  M_n=\sum_{A\in\G_n}M_A,
\end{equation}
where $M_A$ counts the pairs $\{B,2A-B\}$ in $\G_n$, i.e.\ the collinear
``isosceles'' configurations in which $A$ is the midpoint of the base. With
$m_x=\min(x,n-1-x)$, $m_y=\min(y,n-1-y)$,
\[
  M_{(x,y)}=\frac{(2m_x+1)(2m_y+1)-1}{2}.
\]
For even $n$ the values $2m_x+1$ are $1,3,\dots,n-1,n-1,\dots,3,1$, whose sum is
$n^2/2$, so
\begin{equation}\label{eq:M}
  M_n=\frac{(n^2/2)^2-n^2}{2}=\frac{n^2(n^2-4)}{8}.
\end{equation}
(For $n=10$ this gives $1200$, as in the companion note.) Finally $S_A$ and $M_A$
are invariant under the $8$ symmetries of the square, so it suffices to compute
them for the representatives $0\le x\le y\le n/2-1$; a class with $x=y$ has $4$
elements, any other class $8$.

\section{The $20\times20$ grid}

\subsection{Total and collinear triples}
\[
  T_{20}=\binom{400}{3}=10\,586\,800.
\]
Rows and columns contribute $2\cdot20\cdot\binom{20}{3}=2\cdot20\cdot1140=45\,600$
collinear triples. For the oblique ones we group the terms of \eqref{eq:L} by
$g=\gcd(a,b)$; Table~\ref{tab:oblique} lists the contribution of each $g$ for one
slope sign. Their sum is $29\,538$, so
\[
  L_{20}=45\,600+2\cdot29\,538=45\,600+59\,076=104\,676,
  \qquad
  T_{20}-L_{20}=10\,482\,124 .
\]

\subsection{Isosceles triangles}

\paragraph{Midpoint term.} By \eqref{eq:M},
$M_{20}=\dfrac{400\cdot396}{8}=19\,800$.

\paragraph{Equidistant pairs.} There are $55$ symmetry classes of apices
(Appendix~\ref{app:20}). As in the $10\times10$ case it is instructive to look at
the corner $A=(0,0)$. The $\binom{20}{2}=190$ pairs $\{(a,b),(b,a)\}$ with $a\ne b$
are always equidistant from $A$. All other pairs come from squared radii with
two or more essentially different representations as a sum of two squares
inside the grid. In the $10\times10$ grid there were only $4$ such radii
($25,50,65,85$); in the $20\times20$ grid there are $29$ of them
($25,50,65,85,100,125,130,145,169,170,\dots$), and some have \emph{three}
representations, e.g.
\[
  325=1^2+18^2=6^2+17^2=10^2+15^2,
\]
giving $k_A(325)=6$ points on that circle. Altogether they contribute $118$ extra
pairs, so
\[
  S_{(0,0)}=190+118=308,\qquad M_{(0,0)}=0.
\]
At the other extreme, the four central points such as $(9,9)$ have
$S_A=1379$ and $M_A=180$, hence are apices of $1199$ isosceles triangles,
$3.9$ times as many as a corner. Summing the table of Appendix~\ref{app:20}
weighted by orbit sizes,
\[
  S_{20}=268\,092,\qquad M_{20}=19\,800,\qquad
  I_{20}=268\,092-19\,800=248\,292 .
\]

\subsection{Result}
\begin{theorem*}
For three distinct uniformly random points of the $20\times20$ grid, conditioned
on not being collinear,
\[
  p_{20}=\frac{248\,292}{10\,482\,124}=\frac{62\,073}{2\,620\,531}\approx0.02368718 .
\]
\end{theorem*}
Here $\gcd(248\,292,\,10\,482\,124)=4$, $62\,073=3^3\cdot11^2\cdot19$ and
$2\,620\,531$ is prime, so the fraction is in lowest terms. Without redrawing
collinear triples the probability would be $I_{20}/T_{20}\approx0.023453$.

\section{The $30\times30$ grid}

\subsection{Total and collinear triples}
\[
  T_{30}=\binom{900}{3}=121\,095\,300.
\]
Rows and columns contribute $2\cdot30\cdot\binom{30}{3}=2\cdot30\cdot4060=243\,600$.
The oblique contributions (Table~\ref{tab:oblique}) sum to $193\,542$ per slope
sign, so
\[
  L_{30}=243\,600+2\cdot193\,542=243\,600+387\,084=630\,684,
  \qquad
  T_{30}-L_{30}=120\,464\,616 .
\]

\begin{table}[ht]
\centering
\small
\begin{tabular}{rrr@{\qquad}rrr}
\toprule
$g$ & $n=20$ & $n=30$ & $g$ & $n=20$ & $n=30$\\
\midrule
 2 & 5852 & 30\,272 & 16 & 240 & 2940 \\
 3 & 4966 & 26\,334 & 17 & 144 & 2704 \\
 4 & 3840 & 21\,588 & 18 &  68 & 2448 \\
 5 & 3100 & 17\,900 & 19 &  18 & 2178 \\
 6 & 2540 & 14\,400 & 20 &     & 1900 \\
 7 & 1950 & 12\,570 & 21 &     & 1620 \\
 8 & 1680 & 10\,724 & 22 &     & 1344 \\
 9 & 1320 &  9144 & 23 &     & 1078 \\
10 &  900 &  7200 & 24 &     &  828 \\
11 &  810 &  6650 & 25 &     &  600 \\
12 &  704 &  5940 & 26 &     &  400 \\
13 &  588 &  5100 & 27 &     &  234 \\
14 &  468 &  4160 & 28 &     &  108 \\
15 &  350 &  3150 & 29 &     &   28 \\
\midrule
\multicolumn{6}{r}{total over one slope sign:\quad $n=20$: $29\,538$\qquad
$n=30$: $193\,542$}\\
\bottomrule
\end{tabular}
\caption{Oblique collinear triples for one slope sign, grouped by
$g=\gcd(a,b)$: the entry is $\sum_{\gcd(a,b)=g}(n-a)(n-b)(g-1)$ over
$1\le a,b\le n-1$.}
\label{tab:oblique}
\end{table}

\subsection{Isosceles triangles}

\paragraph{Midpoint term.} By \eqref{eq:M},
$M_{30}=\dfrac{900\cdot896}{8}=100\,800$.

\paragraph{Equidistant pairs.} There are now $120$ symmetry classes
(Appendix~\ref{app:30}). The corner has $\binom{30}{2}=435$ symmetric pairs plus
$348$ extra pairs coming from $76$ squared radii with several representations,
so $S_{(0,0)}=783$. The central points such as $(14,14)$ have $S_A=3585$,
$M_A=420$, and are apices of $3165$ isosceles triangles, $4.0$ times as many as a
corner. The weighted sums give
\[
  S_{30}=1\,561\,208,\qquad M_{30}=100\,800,\qquad
  I_{30}=1\,561\,208-100\,800=1\,460\,408 .
\]

\subsection{Result}
\begin{theorem*}
For three distinct uniformly random points of the $30\times30$ grid, conditioned
on not being collinear,
\[
  p_{30}=\frac{1\,460\,408}{120\,464\,616}=\frac{7937}{654\,699}\approx0.01212313 .
\]
\end{theorem*}
Here $\gcd(1\,460\,408,\,120\,464\,616)=184=2^3\cdot23$, $7937$ is prime and
$654\,699=3\cdot218\,233$ with $218\,233$ prime, so the fraction is in lowest
terms. Without redrawing collinear triples the probability would be
$I_{30}/T_{30}\approx0.012060$.

\section{Comparison with the $10\times10$ grid}

\subsection{The numbers side by side}

\begin{center}
\begin{tabular}{lrrr}
\toprule
 & $10\times10$ & $20\times20$ & $30\times30$ \\
\midrule
grid points $n^2$ & $100$ & $400$ & $900$ \\
triples $T_n$ & $161\,700$ & $10\,586\,800$ & $121\,095\,300$ \\
collinear triples $L_n$ & $4448$ & $104\,676$ & $630\,684$ \\
collinear share $L_n/T_n$ & $2.75\%$ & $0.99\%$ & $0.52\%$ \\
triangles $T_n-L_n$ & $157\,252$ & $10\,482\,124$ & $120\,464\,616$ \\
equidistant pairs $S_n$ & $12\,484$ & $268\,092$ & $1\,561\,208$ \\
midpoint pairs $M_n$ & $1200$ & $19\,800$ & $100\,800$ \\
isosceles triangles $I_n$ & $11\,284$ & $248\,292$ & $1\,460\,408$ \\
apex classes & $15$ & $55$ & $120$ \\
corner apex count $S_A-M_A$ & $59$ & $308$ & $783$ \\
central apex count $S_A-M_A$ & $217$ & $1199$ & $3165$ \\
\midrule
$p_n$ (exact) & $\dfrac{2821}{39\,313}$ & $\dfrac{62\,073}{2\,620\,531}$ &
  $\dfrac{7937}{654\,699}$ \\[2ex]
$p_n$ (decimal) & $0.0717574$ & $0.0236872$ & $0.0121231$ \\
about one round in & $13.9$ & $42.2$ & $82.5$ \\
$p_n/p_{10}$ & $1$ & $0.3301$ & $0.1689$ \\
$n^2 p_n$ & $7.176$ & $9.475$ & $10.911$ \\
\bottomrule
\end{tabular}
\end{center}

\subsection{How the counts grow}

Going from $n=10$ to $n=20$ and $n=30$ multiplies the side by $2$ and $3$. The
number of triples grows like $n^6/6$: by factors $65.5$ and $749$. The number of
isosceles triangles grows much more slowly, by factors
\[
  \frac{I_{20}}{I_{10}}\approx22.0,\qquad\frac{I_{30}}{I_{10}}\approx129.4,
\]
which is a little faster than $n^4$ (factors $16$ and $81$). Consequently the
probability falls:
\[
  \frac{p_{20}}{p_{10}}\approx0.330\approx\frac13,
  \qquad
  \frac{p_{30}}{p_{10}}\approx0.169\approx\frac16 .
\]
A pure $1/n^2$ law would predict $1/4$ and $1/9$, a $1/n$ law $1/2$ and $1/3$; the
observed decay lies in between, closer to $1/n^2$, and the slowly increasing
column $n^2p_n$ ($7.18\to9.47\to10.91$) shows the missing factor is logarithmic.

\subsection{A heuristic for the $(\ln n)/n^2$ decay}

Count isosceles triangles by their \emph{base} $\{B,C\}$ instead of their apex: the
apex is a lattice point of the perpendicular bisector of $BC$, other than the
midpoint. Let $v=C-B=(a,b)$ and $g=\gcd(a,b)$. When the bisector contains lattice
points at all (which happens for a positive proportion of displacements), these
are spaced $|v|/g$ apart along it, so a bisector crossing the grid contains about
$n g/|v|$ of them. Each displacement $v$ occurs for about $n^2$ bases, so
\[
  I_n\approx c\,n^3\sum_{0<|v|\lesssim n}\frac{\gcd(v)}{|v|}.
\]
Writing $v=g u$ with $u$ primitive, and using that a proportion $6/\pi^2$ of the
lattice vectors are primitive,
\[
  \sum_{0<|v|\le R}\frac{\gcd(v)}{|v|}
  =\sum_{g\le R}\;\sum_{\substack{u\text{ primitive}\\|u|\le R/g}}\frac1{|u|}
  \approx\sum_{g\le R}\frac{12}{\pi}\cdot\frac Rg
  \approx\frac{12}{\pi}R\ln R .
\]
Hence $I_n$ is of order $n^4\ln n$, and since $T_n-L_n\sim n^6/6$,
\[
  p_n\approx C\,\frac{\ln n}{n^2}.
\]
The $\ln n$ comes from bases with a large common divisor $g$, whose bisectors are
densely populated by lattice points. This is only a heuristic (boundary effects
and the parity condition for the bisector to meet the lattice are absorbed into
the constants), but it matches the exact values well: computing $p_n$ exactly
with the same method for all even $n$ from $10$ to $40$ and for $n=50,60$, a
least-squares fit gives
\[
  n^2p_n\approx3.51\ln n-1.00,
\]
with maximal error $0.10$ over that range. For the three grids of this note
it predicts $n^2p_n\approx7.08,\ 9.51,\ 10.93$, against the exact
$7.18,\ 9.47,\ 10.91$. This fitted slope is still affected by lower-order terms;
Section~\ref{sec:general} makes the heuristic precise and shows that the true
asymptotic slope is $24/\pi^2+4/\pi\approx3.705$.

\subsection{Why larger grids favour the centre more}

In all three grids the corners are the least productive apices and the centre
the most productive, because more circles around a central point stay inside
the grid. The ratio centre/corner grows slowly: $217/59\approx3.7$,
$1199/308\approx3.9$, $3165/783\approx4.0$. Larger grids also contain radii with
more representations as sums of two squares (up to three in the corner of the
$20\times20$ and $30\times30$ grids, against two in the $10\times10$ grid),
which adds isosceles triangles that are not mirror images across a diagonal.

\subsection{Collinear triples matter less}

The collinear share $L_n/T_n$ drops from $2.75\%$ to $0.99\%$ and $0.52\%$, so the
effect of redrawing collinear triples shrinks: the ratio $p_n/(I_n/T_n)=T_n/(T_n-L_n)$
is $1.028$, $1.010$ and $1.005$ respectively.

\section{The general $n\times n$ grid}\label{sec:general}

\subsection{What has a closed form, and what does not}

Two of the four quantities in \eqref{eq:p}--\eqref{eq:I} are polynomials in $n$:
\[
  T_n=\binom{n^2}{3}=\frac{n^2(n^2-1)(n^2-2)}{6},
  \qquad
  M_n=\begin{cases}
    \dfrac{n^2(n^2-4)}{8}, & n \text{ even},\\[1.5ex]
    \dfrac{(n^2-1)^2}{8}, & n \text{ odd}.
  \end{cases}
\]
(For odd $n$ the values $2m_x+1$ are $1,3,\dots,n-2,n,n-2,\dots,1$, with sum
$(n^2+1)/2$.) The other two are not. $L_n$ is given exactly by the gcd sum
\eqref{eq:L}; it is the sequence of collinear triples in an $n\times n$ grid
(OEIS A000938), it grows like a constant times $n^4\ln n$, and the gcd prevents it
from collapsing to a polynomial. The isosceles count $I_n$ depends on how many ways
integers can be written as sums of two squares, as we now show, and therefore also
has no polynomial or rational closed form. Consequently $p_n$ has no closed form,
but it has an exact finite-sum formula and a clean asymptotic law.

\subsection{An exact formula for $I_n$}

For integer vectors $u=(u_1,u_2)$ and $w=(w_1,w_2)$ let
\[
  \Delta_x(u,w)=\max(0,u_1,w_1)-\min(0,u_1,w_1),\qquad
  \Delta_y(u,w)=\max(0,u_2,w_2)-\min(0,u_2,w_2)
\]
be the width and height of the smallest box containing $0$, $u$ and $w$, and write
$t_+=\max(t,0)$.

\begin{proposition}\label{prop:exact}
For every $n\ge1$,
\begin{equation}\label{eq:Iexact}
  I_n=\frac12\sum_{\substack{u,w\in\mathbb{Z}^2\setminus\{0\}\\
      |u|^2=|w|^2,\ w\ne\pm u}}
  \bigl(n-\Delta_x(u,w)\bigr)_+\bigl(n-\Delta_y(u,w)\bigr)_+ .
\end{equation}
\end{proposition}
\begin{proof}
Assign to an isosceles triangle with apex $A$ and other vertices $B,C$ the legs
$u=B-A$ and $w=C-A$; then $|u|=|w|$, and since $B\ne C$ and the triangle is not
degenerate, $w\ne\pm u$. Conversely, a pair $(u,w)$ with $|u|=|w|$ and $w\ne\pm u$
together with a point $A$ such that $A$, $A+u$ and $A+w$ lie in $\G_n$ gives an
isosceles triangle with apex $A$. The number of such $A$ is
$(n-\Delta_x)_+(n-\Delta_y)_+$. By Lemma~2 of the $10\times10$ note the apex is
unique, so each triangle arises from exactly the two ordered pairs $(u,w)$ and
$(w,u)$, which explains the factor $\tfrac12$.
\end{proof}

The sum is finite (only $|u_i|,|w_i|\le n-1$ contribute). Evaluating
\eqref{eq:Iexact} by computer reproduces the counts of \texttt{core/exact.py},
e.g.\ $I_{10}=11\,284$, $I_{13}=36\,812$ and $I_{20}=248\,292$.

Grouping the terms by the common squared length $k=|u|^2=|w|^2$ shows where the
arithmetic enters. The circle $|v|^2=k$ carries
\[
  r_2(k)=\#\{v\in\mathbb{Z}^2 : |v|^2=k\}=4\sum_{d\mid k}\chi(d)
\]
lattice points, where $\chi(d)=0,1,-1$ for $d$ even, $d\equiv1$ and $d\equiv3
\pmod4$ respectively. So the circle contributes $r_2(k)^2-2r_2(k)$ ordered pairs
$(u,w)$ to \eqref{eq:Iexact}. Since $r_2(k)$ jumps irregularly ($r_2(3)=0$,
$r_2(25)=12$, $r_2(325)=24$), no polynomial in $n$ can reproduce $I_n$.

\subsection{The asymptotic formula}

Put $h(x,y)=(1-\Delta_x(x,y))_+(1-\Delta_y(x,y))_+$ for real vectors $x,y$, so that
the summand in \eqref{eq:Iexact} is $n^2h(u/n,w/n)$. For $s>0$ let
\[
  H(s)=\frac1{4\pi^2}\int_0^{2\pi}\!\!\int_0^{2\pi}
  h\bigl(\sqrt s\,e^{i\theta},\sqrt s\,e^{i\varphi}\bigr)\,d\theta\,d\varphi
\]
be the average of $h$ over two independent uniformly random points on the circle
of squared radius $s$, and
\[
  K=\int_0^\infty H(s)\,ds .
\]

\paragraph{Derivation.} Two ingredients are needed:
\begin{enumerate}
  \item (Ramanujan) $\displaystyle\sum_{k\le X}r_2(k)^2=4X\ln X+O(X)$, while
    $\sum_{k\le X}r_2(k)=\pi X+O(\sqrt X)$, so the excluded pairs $w=\pm u$ are of
    lower order.
  \item (Equidistribution) On average over $k$, the pairs $(u,w)$ on the circle
    $|v|^2=k$ are spread uniformly in angle, so that
    $\sum_{u,w:\,|u|^2=|w|^2=k}h(u/n,w/n)\approx r_2(k)^2H(k/n^2)$.
\end{enumerate}
With these,
\[
  I_n\approx\frac{n^2}{2}\sum_k r_2(k)^2H\!\Bigl(\frac k{n^2}\Bigr)
  \approx\frac{n^2}{2}\int_0^\infty H\!\Bigl(\frac x{n^2}\Bigr)\,4\ln x\,dx
  =2n^4\int_0^\infty H(s)\ln(n^2s)\,ds
  =4K\,n^4\ln n+O(n^4).
\]
Since $T_n-L_n=\tfrac16n^6\bigl(1+O(n^{-2}\ln n)\bigr)$, this gives
\begin{equation}\label{eq:asym}
  p_n=\frac{I_n}{T_n-L_n}=24K\,\frac{\ln n}{n^2}+O\!\Bigl(\frac1{n^2}\Bigr).
\end{equation}
The logarithm comes from $\sum r_2(k)^2$, i.e.\ from the circles that carry
unusually many lattice points. This is the same phenomenon as the large-gcd bases
in the heuristic of the previous section, seen from the apex instead of the base.
Ingredient~2 is a standard equidistribution property of lattice points on circles,
but we have not written out a rigorous proof that it holds with the required
uniformity. Formula~\eqref{eq:asym} should therefore be read as a derivation that
is confirmed numerically below, not as a theorem.

\subsection{Computing the constant $K$}

Fix the two angles and put $X=\Delta_x(e^{i\theta},e^{i\varphi})$,
$Y=\Delta_y(e^{i\theta},e^{i\varphi})$, $m=\max(X,Y)$, $m'=\min(X,Y)$. Since
$\Delta_x$ and $\Delta_y$ scale linearly, with $s=\rho^2$,
\[
  \int_0^\infty(1-\rho X)_+(1-\rho Y)_+\,ds
  =\int_0^{1/m}2\rho(1-\rho X)(1-\rho Y)\,d\rho
  =\frac{2m-m'}{6m^3},
\]
so
\[
  K=\frac1{4\pi^2}\int_0^{2\pi}\!\!\int_0^{2\pi}\frac{2m-m'}{6m^3}\,d\theta\,d\varphi .
\]
The integrand is piecewise smooth: $X$ and $Y$ change formula where
$\cos\varphi=0,\pm\cos\theta$ and $\sin\varphi=0,\pm\sin\theta$, and $m$ switches
where $X=Y$. All these curves have explicit equations, for instance
$\sin\psi-\cos\psi=\cos\theta-\sin\theta$ in the third quadrant. Using the $8$
symmetries of the square to restrict $\theta$ to $[0,\pi/4]$ (with an extra split
at $\theta=\arctan\frac12$, where two of these curves cross), and integrating each
smooth piece with tanh--sinh quadrature at $65$ significant digits, gives
\[
  K=0.15437283133963621670017405100056575958251198991906536509483148\ldots
\]
The integer-relation algorithm PSLQ, run against $1,\pi,\sqrt2,\ln2,\ln(1+\sqrt2)$ and
other natural constants, then finds the relation $6\pi^2K=6+\pi$, which holds to
within $10^{-64}$. Hence
\[
  K=\frac{6+\pi}{6\pi^2},\qquad 4K=\frac{2(6+\pi)}{3\pi^2}\approx0.6175,\qquad
  24K=\frac{24}{\pi^2}+\frac4\pi\approx3.70495 .
\]
(The value is identified to $64$ digits, but we have not evaluated the piecewise
integral symbolically.) A Monte Carlo estimate of the same integral,
$K\approx0.1544\pm0.0002$, is consistent with this value.

\begin{theorem*}[asymptotic law, under ingredient 2]
\[
  p_n\sim\Bigl(\frac{24}{\pi^2}+\frac4\pi\Bigr)\frac{\ln n}{n^2}
  \approx3.705\,\frac{\ln n}{n^2},
  \qquad
  I_n\sim\frac{2(6+\pi)}{3\pi^2}\,n^4\ln n .
\]
\end{theorem*}

\subsection{Comparison with the exact values}

Table~\ref{tab:asym} compares the asymptotic law with the exact $p_n$ computed by
\texttt{core/exact.py} up to $n=150$. Write $A=24/\pi^2+4/\pi$. The column ``local
slope'' is $\bigl(n^2p_n-n'^2p_{n'}\bigr)/\ln(n/n')$ between consecutive rows. It
climbs from $3.32$ towards $A=3.705$, which explains why the fit in the previous
section, made over small grids, gave the smaller slope $3.51$. The gap
$n^2p_n-A\ln n$ settles at a constant: fitting $B+c/n$ to the rows with $n\ge40$
gives $B\approx-1.79$. The resulting two-term approximation
\[
  p_n\approx\frac{A\ln n+B}{n^2},\qquad A=\frac{24}{\pi^2}+\frac4\pi,\quad
  B\approx-1.79,
\]
is within $2\%$ of the exact value already for $n=20$ and within $1\%$ from
$n=30$ on. The leading term alone overestimates $p_n$ by about $10\%$ even at
$n=150$, because $\ln n$ grows so slowly.

\begin{table}[ht]
\centering
\small
\begin{tabular}{rrrrrrrr}
\toprule
$n$ & $I_n$ & $p_n$ (exact) & $n^2p_n$ & local slope & $n^2p_n-A\ln n$ &
$A\ln n/n^2$ & $(A\ln n+B)/n^2$ \\
\midrule
10 & 11\,284 & 0.0717574 & 7.176 & --- & $-1.355$ & 0.0853096 & 0.0674096 \\
20 & 248\,292 & 0.0236872 & 9.475 & 3.317 & $-1.624$ & 0.0277476 & 0.0232726 \\
30 & 1\,460\,408 & 0.0121231 & 10.911 & 3.541 & $-1.690$ & 0.0140014 & 0.0120125 \\
40 & 5\,071\,344 & 0.0074670 & 11.947 & 3.602 & $-1.720$ & 0.0085419 & 0.0074232 \\
60 & 28\,925\,132 & 0.0037291 & 13.425 & 3.644 & $-1.745$ & 0.0042137 & 0.0037165 \\
80 & 98\,698\,384 & 0.0022624 & 14.479 & 3.666 & $-1.756$ & 0.0025367 & 0.0022571 \\
100 & 254\,780\,204 & 0.0015302 & 15.302 & 3.687 & $-1.760$ & 0.0017062 & 0.0015272 \\
150 & 1\,416\,699\,480 & 0.0007466 & 16.798 & 3.691 & $-1.766$ & 0.0008251 & 0.0007455 \\
\bottomrule
\end{tabular}
\caption{Exact probabilities against the asymptotic law, with
$A=24/\pi^2+4/\pi\approx3.70495$ and $B=-1.79$.}
\label{tab:asym}
\end{table}

\section{Comparison with the simulations}

The Monte Carlo experiment (\texttt{grid\_size\_experiment.py}, results in
\texttt{results/grid\_size\_comparison.csv}) used the same rules for all grids:
$10^4$ independent games of $10^6$ rounds, collinear draws redrawn.
\begin{center}
\begin{tabular}{lrrr}
\toprule
 & $10\times10$ & $20\times20$ & $30\times30$ \\
\midrule
exact $p_n$ & $0.0717574$ & $0.0236872$ & $0.0121231$ \\
mean of the $10^4$ game estimates & $0.0717576$ & $0.0236892$ & $0.0121232$ \\
standard error of that mean & $2.5\times10^{-6}$ & $1.5\times10^{-6}$ &
  $1.1\times10^{-6}$ \\
deviation in standard errors ($z$) & $0.06$ & $1.29$ & $0.11$ \\
std.\ across games (observed) & $2.54\times10^{-4}$ & $1.53\times10^{-4}$ &
  $1.10\times10^{-4}$ \\
std.\ across games (binomial theory) & $2.58\times10^{-4}$ & $1.52\times10^{-4}$ &
  $1.09\times10^{-4}$ \\
\bottomrule
\end{tabular}
\end{center}
All three simulated means lie within $1.3$ standard errors of the exact values,
and the spread across games matches the binomial prediction
$\sqrt{p_n(1-p_n)/10^6}$ in each case.

\paragraph{Independent check.} Enumerating all $10\,586\,800$ triples of the
$20\times20$ grid and all $121\,095\,300$ triples of the $30\times30$ grid, and
classifying each with exact integer arithmetic, gives the same counts
$L_{20}=104\,676$, $I_{20}=248\,292$, $L_{30}=630\,684$ and $I_{30}=1\,460\,408$.
The values above can be reproduced with
\texttt{core.exact.exact\_isosceles\_probability(n, n)}.

\appendix

\section{Apex classes of the $20\times20$ grid}\label{app:20}

For each representative $(x,y)$ with $0\le x\le y\le9$: $S_A$, $M_A$ and the number
$S_A-M_A$ of isosceles triangles with apex $A$. Classes with $x=y$ have $4$
elements, the others $8$. Weighted by orbit size the columns sum to
$S_{20}=268\,092$, $M_{20}=19\,800$ and $I_{20}=248\,292$.

\begin{center}
\small
\begin{tabular}{crrr@{\quad}crrr@{\quad}crrr}
\toprule
$(x,y)$ & $S_A$ & $M_A$ & $I_A$ & $(x,y)$ & $S_A$ & $M_A$ & $I_A$ &
$(x,y)$ & $S_A$ & $M_A$ & $I_A$\\
\midrule
$(0,0)$ & 308 & 0 & 308 & $(2,2)$ & 470 & 12 & 458 & $(4,8)$ & 807 & 76 & 731 \\
$(0,1)$ & 330 & 1 & 329 & $(2,3)$ & 490 & 17 & 473 & $(4,9)$ & 831 & 85 & 746 \\
$(0,2)$ & 352 & 2 & 350 & $(2,4)$ & 515 & 22 & 493 & $(5,5)$ & 807 & 60 & 747 \\
$(0,3)$ & 367 & 3 & 364 & $(2,5)$ & 550 & 27 & 523 & $(5,6)$ & 834 & 71 & 763 \\
$(0,4)$ & 390 & 4 & 386 & $(2,6)$ & 574 & 32 & 542 & $(5,7)$ & 875 & 82 & 793 \\
$(0,5)$ & 417 & 5 & 412 & $(2,7)$ & 600 & 37 & 563 & $(5,8)$ & 916 & 93 & 823 \\
$(0,6)$ & 435 & 6 & 429 & $(2,8)$ & 632 & 42 & 590 & $(5,9)$ & 937 & 104 & 833 \\
$(0,7)$ & 453 & 7 & 446 & $(2,9)$ & 651 & 47 & 604 & $(6,6)$ & 916 & 84 & 832 \\
$(0,8)$ & 477 & 8 & 469 & $(3,3)$ & 557 & 24 & 533 & $(6,7)$ & 965 & 97 & 868 \\
$(0,9)$ & 489 & 9 & 480 & $(3,4)$ & 584 & 31 & 553 & $(6,8)$ & 1010 & 110 & 900 \\
$(1,1)$ & 391 & 4 & 387 & $(3,5)$ & 625 & 38 & 587 & $(6,9)$ & 1039 & 123 & 916 \\
$(1,2)$ & 409 & 7 & 402 & $(3,6)$ & 648 & 45 & 603 & $(7,7)$ & 1073 & 112 & 961 \\
$(1,3)$ & 425 & 10 & 415 & $(3,7)$ & 678 & 52 & 626 & $(7,8)$ & 1124 & 127 & 997 \\
$(1,4)$ & 452 & 13 & 439 & $(3,8)$ & 711 & 59 & 652 & $(7,9)$ & 1163 & 142 & 1021 \\
$(1,5)$ & 482 & 16 & 466 & $(3,9)$ & 731 & 66 & 665 & $(8,8)$ & 1238 & 144 & 1094 \\
$(1,6)$ & 504 & 19 & 485 & $(4,4)$ & 662 & 40 & 622 & $(8,9)$ & 1277 & 161 & 1116 \\
$(1,7)$ & 527 & 22 & 505 & $(4,5)$ & 707 & 49 & 658 & $(9,9)$ & 1379 & 180 & 1199 \\
$(1,8)$ & 555 & 25 & 530 & $(4,6)$ & 734 & 58 & 676 & & & & \\
$(1,9)$ & 571 & 28 & 543 & $(4,7)$ & 768 & 67 & 701 & & & & \\
\bottomrule
\end{tabular}
\end{center}

\clearpage
\section{Apex classes of the $30\times30$ grid}\label{app:30}

Same layout, representatives $0\le x\le y\le14$. Weighted by orbit size the
columns sum to $S_{30}=1\,561\,208$, $M_{30}=100\,800$ and $I_{30}=1\,460\,408$.

\begin{center}
\small
\begin{tabular}{crrr@{\quad}crrr@{\quad}crrr}
\toprule
$(x,y)$ & $S_A$ & $M_A$ & $I_A$ & $(x,y)$ & $S_A$ & $M_A$ & $I_A$ &
$(x,y)$ & $S_A$ & $M_A$ & $I_A$\\
\midrule
$(0,0)$ & 783 & 0 & 783 & $(2,13)$ & 1489 & 67 & 1422 & $(6,11)$ & 1977 & 149 & 1828 \\
$(0,1)$ & 815 & 1 & 814 & $(2,14)$ & 1521 & 72 & 1449 & $(6,12)$ & 2017 & 162 & 1855 \\
$(0,2)$ & 855 & 2 & 853 & $(3,3)$ & 1182 & 24 & 1158 & $(6,13)$ & 2070 & 175 & 1895 \\
$(0,3)$ & 890 & 3 & 887 & $(3,4)$ & 1219 & 31 & 1188 & $(6,14)$ & 2118 & 188 & 1930 \\
$(0,4)$ & 921 & 4 & 917 & $(3,5)$ & 1265 & 38 & 1227 & $(7,7)$ & 1874 & 112 & 1762 \\
$(0,5)$ & 957 & 5 & 952 & $(3,6)$ & 1306 & 45 & 1261 & $(7,8)$ & 1943 & 127 & 1816 \\
$(0,6)$ & 993 & 6 & 987 & $(3,7)$ & 1352 & 52 & 1300 & $(7,9)$ & 2016 & 142 & 1874 \\
$(0,7)$ & 1027 & 7 & 1020 & $(3,8)$ & 1400 & 59 & 1341 & $(7,10)$ & 2075 & 157 & 1918 \\
$(0,8)$ & 1066 & 8 & 1058 & $(3,9)$ & 1451 & 66 & 1385 & $(7,11)$ & 2146 & 172 & 1974 \\
$(0,9)$ & 1103 & 9 & 1094 & $(3,10)$ & 1500 & 73 & 1427 & $(7,12)$ & 2188 & 187 & 2001 \\
$(0,10)$ & 1139 & 10 & 1129 & $(3,11)$ & 1549 & 80 & 1469 & $(7,13)$ & 2246 & 202 & 2044 \\
$(0,11)$ & 1176 & 11 & 1165 & $(3,12)$ & 1583 & 87 & 1496 & $(7,14)$ & 2296 & 217 & 2079 \\
$(0,12)$ & 1206 & 12 & 1194 & $(3,13)$ & 1623 & 94 & 1529 & $(8,8)$ & 2099 & 144 & 1955 \\
$(0,13)$ & 1238 & 13 & 1225 & $(3,14)$ & 1657 & 101 & 1556 & $(8,9)$ & 2176 & 161 & 2015 \\
$(0,14)$ & 1263 & 14 & 1249 & $(4,4)$ & 1327 & 40 & 1287 & $(8,10)$ & 2237 & 178 & 2059 \\
$(1,1)$ & 906 & 4 & 902 & $(4,5)$ & 1378 & 49 & 1329 & $(8,11)$ & 2310 & 195 & 2115 \\
$(1,2)$ & 943 & 7 & 936 & $(4,6)$ & 1422 & 58 & 1364 & $(8,12)$ & 2362 & 212 & 2150 \\
$(1,3)$ & 977 & 10 & 967 & $(4,7)$ & 1472 & 67 & 1405 & $(8,13)$ & 2425 & 229 & 2196 \\
$(1,4)$ & 1014 & 13 & 1001 & $(4,8)$ & 1528 & 76 & 1452 & $(8,14)$ & 2473 & 246 & 2227 \\
$(1,5)$ & 1056 & 16 & 1040 & $(4,9)$ & 1581 & 85 & 1496 & $(9,9)$ & 2340 & 180 & 2160 \\
$(1,6)$ & 1095 & 19 & 1076 & $(4,10)$ & 1629 & 94 & 1535 & $(9,10)$ & 2401 & 199 & 2202 \\
$(1,7)$ & 1133 & 22 & 1111 & $(4,11)$ & 1682 & 103 & 1579 & $(9,11)$ & 2480 & 218 & 2262 \\
$(1,8)$ & 1178 & 25 & 1153 & $(4,12)$ & 1718 & 112 & 1606 & $(9,12)$ & 2540 & 237 & 2303 \\
$(1,9)$ & 1218 & 28 & 1190 & $(4,13)$ & 1761 & 121 & 1640 & $(9,13)$ & 2607 & 256 & 2351 \\
$(1,10)$ & 1258 & 31 & 1227 & $(4,14)$ & 1795 & 130 & 1665 & $(9,14)$ & 2657 & 275 & 2382 \\
$(1,11)$ & 1300 & 34 & 1266 & $(5,5)$ & 1500 & 60 & 1440 & $(10,10)$ & 2561 & 220 & 2341 \\
$(1,12)$ & 1329 & 37 & 1292 & $(5,6)$ & 1547 & 71 & 1476 & $(10,11)$ & 2654 & 241 & 2413 \\
$(1,13)$ & 1362 & 40 & 1322 & $(5,7)$ & 1607 & 82 & 1525 & $(10,12)$ & 2713 & 262 & 2451 \\
$(1,14)$ & 1391 & 43 & 1348 & $(5,8)$ & 1664 & 93 & 1571 & $(10,13)$ & 2778 & 283 & 2495 \\
$(2,2)$ & 1043 & 12 & 1031 & $(5,9)$ & 1720 & 104 & 1616 & $(10,14)$ & 2843 & 304 & 2539 \\
$(2,3)$ & 1079 & 17 & 1062 & $(5,10)$ & 1770 & 115 & 1655 & $(11,11)$ & 2838 & 264 & 2574 \\
$(2,4)$ & 1116 & 22 & 1094 & $(5,11)$ & 1830 & 126 & 1704 & $(11,12)$ & 2903 & 287 & 2616 \\
$(2,5)$ & 1159 & 27 & 1132 & $(5,12)$ & 1870 & 137 & 1733 & $(11,13)$ & 2980 & 310 & 2670 \\
$(2,6)$ & 1199 & 32 & 1167 & $(5,13)$ & 1922 & 148 & 1774 & $(11,14)$ & 3049 & 333 & 2716 \\
$(2,7)$ & 1240 & 37 & 1203 & $(5,14)$ & 1960 & 159 & 1801 & $(12,12)$ & 3067 & 312 & 2755 \\
$(2,8)$ & 1287 & 42 & 1245 & $(6,6)$ & 1673 & 84 & 1589 & $(12,13)$ & 3152 & 337 & 2815 \\
$(2,9)$ & 1330 & 47 & 1283 & $(6,7)$ & 1735 & 97 & 1638 & $(12,14)$ & 3217 & 362 & 2855 \\
$(2,10)$ & 1377 & 52 & 1325 & $(6,8)$ & 1794 & 110 & 1684 & $(13,13)$ & 3332 & 364 & 2968 \\
$(2,11)$ & 1423 & 57 & 1366 & $(6,9)$ & 1859 & 123 & 1736 & $(13,14)$ & 3407 & 391 & 3016 \\
$(2,12)$ & 1452 & 62 & 1390 & $(6,10)$ & 1916 & 136 & 1780 & $(14,14)$ & 3585 & 420 & 3165 \\
\bottomrule
\end{tabular}
\end{center}

\end{document}
